\section*{Streszczenie}
\addcontentsline{toc}{section}{Streszczenie}

Coraz więcej danych trafia do chmury obliczeniowej, a wraz z nimi rośnie ryzyko
ich naruszenia. Celem tej pracy jest sprawdzenie, jak dobrze metody uczenia
maszynowego potrafią wykrywać ataki prowadzące do takich naruszeń, oraz
porównanie dwóch popularnych podejść: klasycznego \emph{lasu losowego} oraz
\emph{sieci głębokich} (MLP, CNN i LSTM). Zadanie sformułowano jako rozpoznawanie
złośliwych połączeń sieciowych --- model decyduje, czy dane połączenie to ruch
normalny, czy atak.

W ramach projektu zbudowano kompletny, samodzielnie uruchamialny program w
języku Python, który wczytuje dane, przygotowuje je, trenuje wszystkie modele i
porównuje ich wyniki. Eksperyment przeprowadzono na danych \texttt{\resDataSource},
opartych na zbiorze CSE-CIC-IDS2018 --- ruchu sieciowym zebranym na
infrastrukturze chmurowej. Modele oceniono zarówno pod kątem skuteczności
(m.in.\ wykrywalność ataków i liczba fałszywych alarmów), jak i kosztu (czas
uczenia oraz złożoność).

Wyniki pokazują, że najlepszy model głęboki (\resBestDeepModel,
F1\,=\,\resBestDeepFOne) osiąga nieco wyższą skuteczność niż las losowy
(F1\,=\,\resRFFOne), ale kosztem znacznie dłuższego uczenia. Las losowy okazuje
się najlepszym kompromisem między jakością a kosztem, a bardziej skomplikowane
sieci (CNN, LSTM) nie dają tutaj wyraźnej przewagi.

\medskip
\noindent\textbf{Słowa kluczowe:} bezpieczeństwo chmury, naruszenie danych,
wykrywanie ataków, uczenie maszynowe, las losowy, sieci głębokie,
CSE-CIC-IDS2018.
